% TODO make hw stylesheet
\documentclass[11pt]{scrartcl}
\usepackage[a4paper, total={6in, 8in}]{geometry}
\usepackage[usenames,dvipsnames,svgnames]{xcolor}
\usepackage[shortlabels]{enumitem}
\usepackage[framemethod=TikZ]{mdframed}
\usepackage{amsmath,amssymb,amsthm}
\usepackage[colorlinks]{hyperref}
\usepackage{multicol}
\usepackage{tabularx}

\hypersetup{
	colorlinks=true,
	linkcolor=blue,
	filecolor=magenta,      
	urlcolor=magenta,
	pdftitle={MAT 320 HW}
}

\usepackage{microtype}
\usepackage{mathtools}
\usepackage[headsepline]{scrlayer-scrpage}
\usepackage{thmtools}
\usepackage[textsize=footnotesize]{todonotes}

\mdfdefinestyle{mdbluebox}{roundcorner=10pt,innerbottommargin=9pt,
	linecolor=blue,backgroundcolor=TealBlue!5,}
\declaretheoremstyle[headfont=\sffamily\bfseries\color{MidnightBlue},
mdframed={style=mdbluebox},]{thmbluebox}

\mdfdefinestyle{mdbloobox}{frametitlefont=\bfseries,innerbottommargin=8pt,
	nobreak=true,backgroundcolor=TealBlue!5,linecolor=blue,}
\declaretheoremstyle[headfont=\bfseries\color{MidnightBlue},
mdframed={style=mdbloobox},headpunct={\\[3pt]},postheadspace=0pt,]{thmbloobox}

\mdfdefinestyle{mdredbox}{frametitlefont=\bfseries,innerbottommargin=8pt,
	nobreak=true,backgroundcolor=Salmon!5,linecolor=RawSienna,}
\declaretheoremstyle[headfont=\bfseries\color{RawSienna},
mdframed={style=mdredbox},headpunct={\\[3pt]},postheadspace=0pt,]{thmredbox}

\mdfdefinestyle{mdgreenbox}{linecolor=ForestGreen,backgroundcolor=ForestGreen!5,
	linewidth=2pt,rightline=false,leftline=true,topline=false,bottomline=false,}
\declaretheoremstyle[headfont=\bfseries\sffamily\color{ForestGreen!70!black},
mdframed={style=mdgreenbox},headpunct={ --- },]{thmgreenbox}

\mdfdefinestyle{mdorangebox}{linecolor=RedOrange,backgroundcolor=RedOrange!5,
	linewidth=2pt,rightline=false,leftline=true,topline=false,bottomline=false,}
\declaretheoremstyle[headfont=\bfseries\sffamily\color{RedOrange!70!black},
mdframed={style=mdorangebox},headpunct={ --- },]{thmorangebox}

\mdfdefinestyle{mdblackbox}{linecolor=black,backgroundcolor=RedViolet!5!gray!5,
	linewidth=3pt,rightline=false,leftline=true,topline=false,bottomline=false,}
\declaretheoremstyle[mdframed={style=mdblackbox}]{thmblackbox}

\declaretheoremstyle[spaceabove=6pt,spacebelow=6pt]{basehead}

\declaretheorem[style=basehead,name=Problem]{problem}
\declaretheorem[style=thmredbox,name=Problem,numbered=no]{problem*}
\declaretheorem[style=thmbloobox,name=Setup]{setup} % for config geo
\declaretheorem[style=thmredbox,name=Example,sibling=problem]{example}
\declaretheorem[style=thmredbox,name=$\bigstar$ Problem,sibling=problem]{niceprob}
\declaretheorem[style=thmbluebox,name=Theorem,sibling=problem]{theorem}
\declaretheorem[style=thmbluebox,name=Corollary,sibling=theorem]{corollary}
\declaretheorem[style=thmbluebox,name=Proposition,sibling=theorem]{prop}
\declaretheorem[style=thmbluebox,name=Lemma,numberwithin=problem]{lemma}
\declaretheorem[style=thmbluebox,name=Theorem,numbered=no]{theorem*}
\declaretheorem[style=thmbluebox,name=Lemma,numbered=no]{lemma*}
\declaretheorem[style=thmgreenbox,name=Claim,numberwithin=problem]{claim}
\declaretheorem[style=thmgreenbox,name=Claim,numbered=no]{claim*}
\declaretheorem[style=thmblackbox,name=Remark,numberwithin=problem]{remark}
\declaretheorem[style=thmblackbox,name=Remark,numbered=no]{remark*}
\declaretheorem[style=thmgreenbox,name=Definition,sibling=theorem]{definition}
\declaretheorem[style=thmgreenbox,name=Definition,numbered=no]{definition*}
\declaretheorem[style=thmorangebox,name=Warning,sibling=theorem]{warning}
\declaretheorem[style=thmorangebox,name=Warning,numbered=no]{warning*}
\declaretheorem[style=thmorangebox,name=Note,numbered=no]{note}
\declaretheorem[style=thmblackbox,name=Exercise,sibling=problem]{exercise}
\declaretheorem[style=thmblackbox,name=Exercise,numbered=no]{exercise*}
\declaretheorem[style=thmblackbox,name=Question,sibling=problem]{question}
\declaretheorem[style=thmblackbox,name=Question,numbered=no]{question*}

\addtolength{\textheight}{3.14cm}
\providecommand{\ol}{\overline}
\providecommand{\ul}{\underline}
\providecommand{\eps}{\varepsilon}
\providecommand{\half}{\frac{1}{2}}
\providecommand{\dang}{\measuredangle} %% Directed angle
\providecommand{\CC}{\mathbb C}
\providecommand{\FF}{\mathbb F}
\providecommand{\NN}{\mathbb N}
\providecommand{\QQ}{\mathbb Q}
\providecommand{\RR}{\mathbb R}
\providecommand{\ZZ}{\mathbb Z}
\providecommand{\dg}{^\circ}
\providecommand{\ii}{\item}
\providecommand{\setdiff}{\smallsetminus}
\providecommand{\alert}[1]{{\sffamily\textbf{\textcolor{blue}{#1}}}}
\providecommand{\inv}{^{-1}}
\providecommand{\mb}[1]{\mathbf{#1}}

\reversemarginpar

\newenvironment{soln}{\begin{proof}[Solution]}{\end{proof}}
\newenvironment{walkthrough}{\noindent \textbf{Walkthrough.}}{\medskip}
\newlist{walk}{enumerate}{3}
\setlist[walk]{label=\bfseries (\alph*)}

\providecommand{\pow}{\operatorname{Pow}}
\providecommand{\vocab}[1]{{\textbf{\textcolor{ForestGreen}{#1}}}}
\providecommand{\lcm}{\operatorname{lcm}}
\providecommand{\abs}[1]{\left|#1\right|}

\usepackage{asymptote}
\begin{asydef}
	defaultpen(fontsize(10pt));
	unitsize(1cm);
	size(8cm); // set a reasonable default
	usepackage("amsmath");
	usepackage("amssymb");
	settings.tex="pdflatex";
	settings.outformat="pdf";
	// Replacement for olympiad+cse5 which is not standard
	import geometry;
	// recalibrate fill and filldraw for conics
	void filldraw(picture pic = currentpicture, conic g, pen fillpen=defaultpen, pen drawpen=defaultpen)
	{ filldraw(pic, (path) g, fillpen, drawpen); }
	void fill(picture pic = currentpicture, conic g, pen p=defaultpen)
	{ filldraw(pic, (path) g, p); }
	// some geometry
	pair foot(pair P, pair A, pair B) { return foot(triangle(A,B,P).VC); }
	pair orthocenter(pair A, pair B, pair C) { return orthocentercenter(A,B,C); }
	pair centroid(pair A, pair B, pair C) { return (A+B+C)/3; }
	// cse5 abbreviations
	path CP(pair P, pair A) { return circle(P, abs(A-P)); }
	path CR(pair P, real r) { return circle(P, r); }
	pair IP(path p, path q) { return intersectionpoints(p,q)[0]; }
	pair OP(path p, path q) { return intersectionpoints(p,q)[1]; }
	path Line(pair A, pair B, real a=0.6, real b=a) { return (a*(A-B)+A)--(b*(B-A)+B); }
	// cse5 more useful functions
	picture CC() {
		picture p=rotate(0)*currentpicture;
		currentpicture.erase();
		return p;
	}
	pair MP(Label s, pair A, pair B = plain.S, pen p = defaultpen) {
		Label L = s;
		L.s = "$"+s.s+"$";
		label(L, A, B, p);
		return A;
	}
	pair Drawing(Label s = "", pair A, pair B = plain.S, pen p = defaultpen) {
		dot(MP(s, A, B, p), p);
		return A;
	}
	path Drawing(path g, pen p = defaultpen, arrowbar ar = None) {
		draw(g, p, ar);
		return g;
	}
\end{asydef}

\usepackage{mathrsfs}

\ihead{\footnotesize MAT 320 HW09}
\ohead{\footnotesize November 8, 2024}

\title{\vspace{-2em}MAT 320 HW09}
\author{\large Aditya Pahuja}
\date{\large Due 11/8}
\begin{document}
\maketitle

\begin{problem}
	Show that $f\colon(1,+\infty)\to\RR$ defined by $f(x) = \log x$
	is uniformly continuous.
\end{problem}

\begin{soln}
	We first prove the following inequality:
	
	\begin{claim}
		For all $x\ge y > 1$,
		\[\frac{e^y}{y} \le \frac{e^x}{x}.\]
	\end{claim}
	\begin{proof}
		Since $\frac{1}{y} - \frac1x = \frac{x-y}{xy} \le x - y$,
		\begin{align*}
			\frac{1}{y} - \frac{1}{x} &\le x - y \le (e-1)(x - y)\\
			&= \sum_{k=1}^\infty \frac{x-y}{k!}
			= \sum_{k=2}^\infty \frac{(x - y)k}{k!}\\
			&\le \sum_{k=2}^\infty
			\frac{(x - y)(x^{k-2} + x^{k-3}y + \cdots + xy^{k-3} + y^{k-2})}{k!}
			\\ &= \sum_{k=2}^\infty \frac{x^{k-1} - y^{k-1}}{k!}\\
			&= \left(\frac{e^x}{x} - 1 - \frac 1x\right) -
			\left(\frac{e^y}{y} - 1 - \frac 1y\right),
		\end{align*}
		which simplifies directly to the inequality we want.
	\end{proof}
	
	Now, $\log$ is increasing
	(because its inverse $\exp$ is also increasing),
	so we can take logs to get
	\[y - \log y \le x - \log x
	\iff \log x - \log y \le x - y.\]
	This implies uniform continuity,
	since $x - y < \eps$ implies $\log x - \log y < \eps$
	for all $\eps > 0$.
\end{soln}

\pagebreak
\begin{problem}
	Let $(X, d)$ be a metric space and $E$ be a subset of $X$. Define
	\[f_E(x) = \inf_{y\in E} d(x, y).\]
	Show that
	\begin{enumerate}[(a)]
		\ii $f_E$ is uniformly continuous.
		\ii $f_E(x) = 0$ if and only if $x\in\ol E$.
		\ii if $A$ and $B$ are disjoint closed sets then
		\[f(x) = \frac{f_A(x)}{f_A(x) + f_B(x)}\]
		is a continuous function such that $f = 0$
		precisely on the set $A$ and $f = 1$ precisely on the set $B$.
	\end{enumerate}
\end{problem}

\begin{soln}
	Let $a$, $b$ be in $X$ and let $\eps > 0$.
	Suppose $d(a, b) < \eps$.
	\begin{itemize}
		\ii If one of $a$ and $b$ is in $E$ (say WLOG $b$), then
		$f_E(b) = 0$ because it is nonnegative and bounded above by
		$d(b, b) = 0$, so
		\[|f_E(a) - f_E(b)| = f_E(a) =
		\inf_{y\in E} d(a, y) \le d(a, b) < \eps,\]
		where the first inequality comes from $b\in E$.
		\ii If neither $a$ nor $b$ is in $E$, then let $\delta > 0$.
		Then, we can pick $x$ such that
		\[d(b, x) \le f_E(b) + \delta.\]
		This lets us bound
		\begin{align*}
			f_E(a) - f_E(b) &= (f_E(a) + \delta) - (f_E(b) + \delta)\\
			&\le d(a, x) - d(b, x) + \delta\\
			&\le d(a, b) + \delta.
		\end{align*}
		This same bound holds for $f_E(b) - f_E(a)$ by symmetry,
		so, taking $\delta$ arbitrarily small, we can conclude that
		\[|f_E(a) - f_E(b)|\le d(a, b) < \eps.\]
	\end{itemize}
	Thus, for any $a$ and $b$,
	\[d(a, b) < \eps \implies |f_E(a) - f_E(b)| < \eps,\]
	so $f_E$ is uniformly continuous, proving part (a).\\
	
	For part (b), suppose first that $x$ is a limit point of $E$
	(so it is in $\ol E$). Then, for each $\eps > 0$,
	there is a $y\in E$ such that $d(x, y) < \eps$;
	by definition of $\inf$, this means that
	\[f_E(x) = \inf_{y\in E} d(x, y) = 0.\]
	Conversely, if $x$ isn't a limit point of $E$, then
	some open ball $B_X(x, \eps)$ contains no points of $E$;
	that is, every point $y$ of $E$ satisfies
	$d(x, y) \ge \eps$, hence
	\[f_E(x) = \inf_{y\in E} d(x, y) \ge \eps > 0.\]
	
	For part (c), we note that $A$ and $B$ being closed
	means $A = \ol A$ and $B = \ol B$.
	Since the sets are disjoint, there is no point at which
	$f_A(x) = f_B(x) = 0$, ensuring that $f$ is well-defined.
	$f$ must then be continuous
	because sums and quotients of continuous functions are continuous
	as long as the divisor of the quotient is nonzero.	
	Moreover, since $f_A$ and $f_B$ are always nonnegative,
	\[0 \le \frac{f_A(x)}{f_A(x) + f_B(x)} \le 1.\]
	The upper bound is achieved if and only if $f_B(x) = 0$,
	i.e. when $x\in B$, and the lower bound is achieved
	if and only if $f_A(x) = 0$, i.e. when $x\in A$.
\end{soln}

\begin{problem}
	Let $(X, d)$ be a metric space and $D\subseteq X$ a dense subset.
	Show that if $f\colon D\to \RR$ is uniformly continuous
	then there exists a unique continuous function
	$g\colon X\to\RR$ such that $g(x) = f(x)$ for all $x\in D$.
\end{problem}

\begin{soln}
	First, we show that such a function $g$ exists.
	\begin{claim}
		Let $(x_n)$ be a sequence of points in $D$ that converges in $X$.
		Then, the sequence $(f(x_n))$ converges in $\RR$.
	\end{claim}
	\begin{proof}
		Note that $(x_n)$ being a convergent sequence means that it is Cauchy,
		as shown in Problem 6.
		By uniform continuity, we know that, for each $\eps > 0$,
		there exists a $\delta > 0$ such that
		\[d(x_m, x_n) < \delta\implies |f(x_m) - f(x_n)| < \eps.\]
		Since $(x_n)$ is Cauchy, there exists $N$ such that
		$m,\,n > N$ implies that $d(x_m, x_n) < \delta$, so
		\[m,\,n > N \implies d(x_m, x_n) < \delta
		\implies |f(x_m) - f(x_n)| < \eps,\]
		which tells us that $(f(x_n))$ is a Cauchy sequence,
		so it converges because $\RR$ is complete.
	\end{proof}

	\begin{claim}
		Let $(Y, d_Y)$ be a metric space.
		If $(a_n)$ is a sequence of points in $Y$ with limit $\ell$ and
		$(b_n)$ is a sequence of points in $Y$ such that
		$d_Y(a_n, b_n) \to 0$, then
		\[\lim_{n\to\infty} b_n = \ell.\]
	\end{claim}
	\begin{proof}
		For each $\eps > 0$, there exists $N$ such that
		\[n > N \implies d_Y(\ell, a_n) < \frac\eps2,\quad
		d_Y(a_n, b_n) < \frac\eps2,\]
		so, for $n > N$,
		\[d_Y(\ell, b_n) \le d_Y(\ell, a_n) + d_Y(a_n, b_n)
		< \frac\eps2 + \frac\eps2 = \eps,\]
		which is what it means for $b_n$ to converge to $\ell$.
	\end{proof}
	
	\begin{claim}
		Let $(x_n)$ and $(y_n)$ be two sequences of points in $D$
		that converge to the same point $\ell\in X$.
		Then,
		\[\lim_{n\to\infty} f(x_n) = \lim_{n\to\infty} f(y_n).\]
	\end{claim}
	\begin{proof}
		Let $\eps > 0$. We know that there exists $\delta > 0$ such that
		\[d(x_n, y_n) < \delta \implies |f(x_n) - f(y_n)| < \eps\]
		because $f$ is uniformly continuous.
		Moreover, since $x_n\to \ell$ and $y_n\to\ell$,
		\[d(x_n, \ell) < \frac\delta2,\qquad d(y_n,\ell) < \frac\delta2\]
		are true eventually, say when $n > N$, so
		\[d(x_n, y_n) < \delta\]
		is true for all $n > N$. That is,
		\[n > N \implies d(x_n, y_n) < \delta \implies
		|f(x_n) - f(y_n)| < \eps.\]
		Thus,
		\[\lim_{n\to\infty} f(x_n) - f(y_n) = 0 \implies
		\lim_{n\to\infty} f(x_n) = \lim_{n\to\infty} f(y_n)\]
		with the implication coming from Claim 3.2
		(the limits both exist by Claim 3.1).
	\end{proof}
	
	Therefore, for any point $x\in X$, we can define
	\[g(x) := \lim_{n\to\infty} f(x_n)\]
	where $(x_n)$ is any sequence of points in $D$ that converges to $x$.
	Such a sequence $(x_n)$ exists because we can take
	$x_n$ to be any point in $D\cap B(x, \frac 1n)$, which is nonempty
	because $D$ is dense in $X$.
	Also, $g(x)$ well-defined because Claim 3.1 says that the limit exists
	and Claim 3.3 says that the value of the limit is only dependent on $x$.
	
	We now have to show that $g(x) = f(x)$ for all $x\in D$
	and $g$ is continuous. For the first part, we just observe that
	if $x\in D$ then we can take $x_n = x$ for all $n$ so that
	\[g(x) = \lim_{n\to\infty} f(x_n) = \lim_{n\to\infty} f(x) = f(x).\]
	
	To show that $g$ is continuous at any $c\in X$,
	let $(x_n)$ be a sequence of points in $X$ converging to $c$.
	By density we can pick a sequence $(y_n)$ of points in $D$
	such that $d(x_n, y_n) < \frac 1n$.
	Since
	\[d(y_n, c) \le d(y_n, x_n) + d(x_n, c)\]
	and the right-hand side has limit zero, we can say that
	$d(y_n, c) \to 0$, implying $y_n\to c$.
	Letting $\eps > 0$,
	uniform continuity says that there exists a $\delta$ such that
	\[d(x_n, y_n) < \delta \implies |f(x_n) - f(y_n)| < \eps,\]
	so the bound on $|f(x_n) - f(y_n)|$ holds as long as $n > \frac 1\delta$.
	That is, for any $\eps > 0$, we have $|f(x_n) - f(y_n)| < \eps$ eventually,
	which lets us say that
	\[\lim_{n\to\infty} f(x_n) - f(y_n) = 0
	\implies \lim_{n\to\infty} f(x_n) = \lim_{n\to\infty} f(y_n) = g(c),\]
	again using Claim 3.2 to deduce that the limits are equal.\\
	
	Now, we show that $g$ is unique. Indeed, suppose that $h\colon X\to\RR$
	is another continuous function such that $h(x) = f(x)$ for all $x\in D$.
	Then, $g - h$ is a continuous function that is zero
	on all of $D$. Now, for any $x\in X$, let $x_n\in D$
	be a point such that $d(x, x_n) < \frac 1n$;
	these $x_n$ exist because $D$ is dense in $X$. By continuity,
	\[g(x) - h(x) = \lim_{n\to\infty} g(x_n) - h(x_n)
	= \lim_{n\to\infty} 0 = 0,\]
	meaning $g(x) = h(x)$ for all $x$, i.e. $g$ and $h$ are the same function.
\end{soln}

\begin{problem}
	Let $(X, d)$ be a metric space.
	Show that if each closed and bounded subset of $X$ is compact,
	then $(X, d)$ is complete.
\end{problem}

\begin{soln}
	Let $(x_n)$ be an arbitrary Cauchy sequence in $X$.
	We showed in class that $S = \{x_n\}$ is bounded,
	so the closure of this set is bounded:
	\begin{claim}
		For a set $C$ with finite diameter,
		$\operatorname{diam}C = \operatorname{diam}\ol C$.
	\end{claim}
	\begin{proof}
		Since $C\subseteq\ol C$,
		\[\operatorname{diam} C \le \operatorname{diam}\ol C.\]
		
		To show the opposite inequality, we prove equivalently that
		\[\operatorname{diam}\ol C \le \operatorname{diam} C + \eps\]
		for every $\eps > 0$.
		Indeed, let $x,\,y\in\ol C$. Then, $x$ and $y$ are limit points of $C$,
		so we can find $x',\,y'\in C$ such that
		$d(x,x')$ and $d(y, y')$ are both less than $\frac\eps 2$.
		The triangle inequality then says that
		\[d(x, y) \le d(x, x') + d(x', y') + d(y', y)
		< \frac\eps 2 + \operatorname{diam}C + \frac\eps2
		= \operatorname{diam} C + \eps.\]
		Thus $\operatorname{diam} C + \eps$ is an upper bound
		for $d(x, y)$, meaning it is at least as large as
		$\operatorname{diam} C$, as desired.
	\end{proof}
	
	This means that $\ol S$ is compact because it is a closed and bounded
	subset of $X$.
	Now, $\ol S$ is sequentially compact, so
	$(x_n)$ has a convergent subsequence $(x_{n_k})$ with limit $\ell$.
	I claim that $(x_n)$ converges to $\ell$ as well.
	Let $\eps > 0$ and pick $N$ so that
	\begin{align*}
		n_k > N &\implies |\ell - x_{n_k}| < \frac\eps2\\
		m,\,n > N &\implies |x_m - x_n| < \frac\eps2,
	\end{align*}
	where the latter inequality comes from $(x_n)$ being Cauchy.
	Then, if $n_k$ and $n$ are both bigger than $N$,
	\[|\ell - x_n| \le |\ell - x_{n_k}| + |x_{n_k} - x_n|
	< \frac\eps2 + \frac\eps2 = \eps,\]
	proving that $x_n\to\ell$, i.e. every Cauchy sequence in $X$ converges.
\end{soln}

\begin{problem}
	Let $(X, d)$ be a metric space. Show that
	\begin{enumerate}[(a)]
		\ii if $E\subseteq X$ is complete in the induced metric then
		$E$ is a closed subset of $X$.
		\ii if $X$ is complete, then a subset $E\subseteq X$ is closed
		if and only if it is complete in the induced metric.
	\end{enumerate}
\end{problem}

\begin{soln}
	For (a), let $\ell$ be a limit point of $E$.
	For every positive integer $n$, there exists an $x_n$
	such that $d(\ell, x_n) < \frac 1n$.
	Then, the sequence $(x_n)$ converges to $\ell$ by construction,
	so, by Problem 6, it is Cauchy. By completeness,
	$(x_n)$ must therefore converge to a point in $E$,
	so $\ell$, the limit of $x_n$, is in $E$.
	Thus $E$ is closed because it contains all of its limit points.\\
	
	For (b), we already showed that complete implies closed.
	Conversely, if $E$ is closed, consider a Cauchy sequence
	$(x_n)$ of points in $E$. Since this sequence is also Cauchy in $X$,
	it has a limit $\ell$. However, $\ell$ is then a limit point of $E$,
	so $E$ being closed means that $\ell\in E$.
	Hence, every Cauchy sequence of points in $E$
	has a limit in $E$, i.e. $E$ is complete.
\end{soln}

\begin{problem}
	Show that if $x_n\to\ell$ in a metric space $(X, d)$,
	then $(x_n)$ is a Cauchy sequence.
\end{problem}

\begin{soln}
	Let $\eps > 0$. Convergence to $\ell$ means that
	there exists an $N$ such that
	\[n > N \implies |x_n - \ell| < \frac\eps2.\]
	Then,
	\[m,\,n > N \implies |x_m - x_n| \le
	|x_m - \ell| + |\ell - x_n| < \eps,\]
	so the sequence is Cauchy.
\end{soln}

\begin{problem}
	Let $(X, d_X)$ be a metric space and $S$ a set.
	We say that a function $f\colon S\to X$ is bounded
	if its image $f(S)\subseteq X$ has finite diameter.
	Let
	\[\mathcal B(S, X) = \{f\colon S\to X : f\text{ is bounded}\}.\]
	Show that
	\begin{enumerate}[(a)]
		\ii if $f$ and $g$ are bounded then
		\[\sup_{s\in S} d_X(f(s), g(s)) < +\infty.\]
		\ii the function
		\[d(f, g) = \sup_{s\in S} d_X(f(s), g(s))\]
		is a metric on $\mathcal B(S, X)$.
		\ii if $X$ is complete, then $(\mathcal B(S, X), d)$
		is a complete metric space.
	\end{enumerate}
\end{problem}

\begin{soln}
	(a) Fix $a\in S$. Then, for any $s$,
	\begin{align*}
		d_X(f(s), g(s)) &\le
		d_X(f(s), f(a)) + d_X(f(a), g(a)) + d_X(g(a), g(s))\\
		&\le d_X(f(a), g(a)) + \operatorname{diam}(f(S)) + \operatorname{diam}(g(S)),
	\end{align*}
	so $d_X(f(s), g(s))$ is bounded above by a constant,
	hence its supremum over all $s$ is real.
	
	(b) Moreover, $d$ is a metric.
	First, it is clearly nonnegative because it is the supremum of
	a collection of nonnegative values.
	Second, it is zero iff $f$ and $g$ are the same function because
	\[\sup_{s\in S} d_X(f(s), g(s)) = 0 \implies
	d_X(f(s), g(s)) \le 0 \text{ for all }s\in S\]
	which implies $d_X(f(s), g(s))$ is identically zero
	and so $f(s) = g(s)$ for all $s\in S$, and
	\[d(f, f) = \sup_{s\in S} d_X(f(s), f(s)) = \sup_{s\in S} 0 = 0.\]
	It is also trivial to see that $d(f, g) = d(g, f)$.
	Finally, the triangle inequality holds because
	\begin{align*}
		d(f, h) + d(g, h) &=
		\sup_{s\in S} d_X(f(s), h(s)) + \sup_{s\in S} d_X(g(s), h(s))\\
		&= \sup_{s\in S} [d_X(f(s), h(s)) + d_X(g(s), h(s))]\\
		&\ge \sup_{s\in S} d_X(f(s), g(s)) = d(f, g).
	\end{align*}
	
	(c) Lastly, suppose $X$ is complete. Then, let
	$(f_n)$ be a Cauchy sequence in $\mathcal B(S, X)$.
	We see that, for each $s\in S$, the sequence
	$(f_n(s))$ is Cauchy in $X$ because, for $\eps > 0$,
	if $N$ is chosen such that $m,\,n > N$ implies
	$d(f_m, f_n) < \eps$, then
	\[d_X(f_m(s), f_n(s)) \le d(f_m, f_n) < \eps.\]
	Thus $(f_n(s))$ converges in $X$ because $X$ is complete,
	so $(f_n)$ has a pointwise limit $f$.
	
	I claim that $f$ is also the limit of $(f_n)$ with respect to $d$.
	Let $\eps > 0$ and take $N$ such that $d(f_m, f_n) < \eps/2$
	whenever $m$ and $n$ exceed $N$. Then, for any $s\in S$,
	$m,\,n > N$ implies
	\begin{align*}
		d_X(f_n(s), f(s)) &\le d_X(f_n(s), f_m(s)) + d_X(f_m(s), f(s))\\
		&\le d(f_n, f_m) + d_X(f_m(s), f(s))\\
		& < \eps/2 + d_X(f_m(s), f(s)).
	\end{align*}
	Taking $m\to\infty$, this means that
	\[d_X(f_n(s), f(s)) \le \eps/2 < \eps\]
	for all $s\in S$, so for any $\eps > 0$ there is an $N$
	such that $n > N$ implies that $d(f_n, f)<\eps$,
	which is what it means for $f_n$ to converge to $f$
	with respect to the metric $d$.
\end{soln}

\begin{problem}
	Let $(K, d_K)$ and $(X, d_X)$ be two metric spaces
	and $C^0(K, X)$ be the space of all continuous functions from $K$ to $X$.
	Show that if $K$ is compact and $X$ is complete then
	$C^0(K, X)$ is a complete metric space with metric
	\[d(f, g) = \max_{x\in K} d_X(f(x), g(x)).\]
\end{problem}

\begin{soln}
	First, we show that $d$ is well-defined; that is,
	\[\max_{x\in K} d_X(f(x), g(x))\]
	actually exists for each $f$ and $g$.
	
	\begin{claim}
		For any continuous functions $f,g\colon K\to X$,
		the function $\alpha\colon K\to\RR$ defined by
		\[\alpha(x) = d_X(f(x), g(x))\]
		is continuous.
	\end{claim}
	\begin{proof}
		Let $c$ be an element of $X$ and let $\eps > 0$. Then,
		\begin{align*}
			\alpha(c) &\le d_X(f(c), f(x)) +
			\alpha(x) + d_X(g(x), g(c)),\\
			\alpha(x) &\le d_X(f(c), f(x)) +
			\alpha(c) + d_X(g(x), g(c)).
		\end{align*}
		This means that we can bound
		\[|\alpha(x) - \alpha(c)| \le d_X(f(c), f(x)) + d_X(g(c), g(x)).\]
		By continuity of $f$ and $g$, we can pick $\delta> 0$ so that
		\[d_K(x, c) < \delta \implies
		d_X(f(x), f(c)) < \eps/3,\quad d_X(g(x), g(c)) < \eps/3.\]
		This means that $d_K(x, c) < \delta$ implies
		\[|\alpha(x) - \alpha(c)| \le \frac{2\eps}3 < \eps,\]
		which is what it means for $\alpha$ to be continuous.
	\end{proof}
	
	Now the Weierstrass theorem implies that $\alpha(K)\subseteq\RR$
	is an interval, so it has a maximum.
	
	To show that $d$ is a metric, we observe the following facts:
	\begin{itemize}
		\ii $d(f, g)$ is nonnegative because it is the maximum
		of a collection of nonnegative values.
		\ii $d(f, g) = 0$ if $f = g$ because
		\[d(f, f) = \max_{x\in K} d_X(f(x), g(x)) =
		\max_{x\in K} 0 = 0,\]
		and conversely, if $d(f, g) = 0$, then
		\[0 \le d_X(f(x), g(x)) \le 0\]
		for all $x\in K$, which forces $d_X(f(x), g(x)) = 0 \iff f(x) = g(x)$
		on all of $K$.
		\ii $d(f, g) = d(g, f)$ is trivial, as $\max(a, b) = \max(b, a)$.
		\ii The triangle inequality holds because
		\begin{align*}
			d(f, h) + d(g, h) &=
			\max_{x\in K} d_X(f(x), h(x)) + \max_{x\in K} d_X(g(x), h(x))\\
			&= \max_{x\in K} [d_X(f(x), h(x)) + d_X(g(x), h(x))]\\
			&\ge \max_{x\in K} d_X(f(x), g(x)) = d(f, g).
		\end{align*}
	\end{itemize}
	
	It remains to prove that $(C^0(K, X), d)$ is complete,
	which can be done in the same way as 7(c).
	Let $(f_n)$ be a Cauchy sequence in $\mathcal C^0(K, X)$.
	We see that, for each $x\in K$, the sequence
	$(f_n(x))$ is Cauchy in $X$ because, for $\eps > 0$,
	if $N$ is chosen such that $m,\,n > N$ implies
	$d(f_m, f_n) < \eps$, then
	\[d_X(f_m(x), f_n(x)) \le d(f_m, f_n) < \eps.\]
	Thus $(f_n(x))$ converges in $X$ because $X$ is complete,
	so $(f_n)$ has a pointwise limit $f$.
	
	I claim that $f$ is also the limit of $(f_n)$ with respect to $d$.
	Let $\eps > 0$ and take $N$ such that $d(f_m, f_n) < \eps/2$
	whenever $m$ and $n$ exceed $N$. Then, for any $x\in K$,
	$m,\,n > N$ implies
	\begin{align*}
		d_X(f_n(x), f(x)) &\le d_X(f_n(x), f_m(x)) + d_X(f_m(x), f(x))\\
		&\le d(f_n, f_m) + d_X(f_m(x), f(x))\\
		& < \eps/2 + d_X(f_m(x), f(x)).
	\end{align*}
	Taking $m\to\infty$, this means that
	\[d_X(f_n(x), f(x)) \le \eps/2 < \eps\]
	for all $x\in K$, so for any $\eps > 0$ there is an $N$
	such that $n > N$ implies that $d(f_n, f)<\eps$,
	which is what it means for $f_n$ to converge to $f$
	with respect to the metric $d$.
\end{soln}

\pagebreak
\begin{problem}
	Prove that $\ell^1$ is a complete metric space.
\end{problem}
\begin{soln}
	Let $(\mb x_n)$ be a Cauchy sequence in $\ell^1$.
	Then, for each $\eps > 0$, there exists $N$ such that
	$m,\,n > N$ implies
	\[\sum_{k=1}^\infty |x_{m, k} - x_{n, k}| < \eps.\]
	Since $|x_{m, k} - x_{n, k}| < \|\mb x_m - \mb x_n\|$
	(where the norm is $\|\mb x\| = d(\mb x, \mb 0)$),
	this means the sequence $(x_{n,k})_{n\in\NN}$
	of the $k$th coordinates of each $\mb x_n$ is Cauchy,
	so, since $\RR$ is complete, $(\mb x_n)$ converges coordinatewise
	(i.e. $(x_{n,k})_n$ converges for every $k$).
	Let $\ell_k$ be the limit of $x_{n,k}$ as $n\to\infty$.
	
	\begin{claim}
		The vector
		\[\mb x = (\ell_1, \ell_2, \ell_3, \dots)\]
		is in $\ell^1$.
	\end{claim}
	\begin{proof}
		We wish to show that
		\[\sum_{i=1}^\infty |\ell_i|\]
		converges. It suffices to show that
		the partial sums are bounded because the numbers being summed
		are all nonnegative.
		
		For each $i\le N$, there exists an $M_i$ such that
		\[n > M_i\implies |\ell_i - x_{n,i}| < 2^{-i}.\]
		Fix $n > \max\{M_1, \dots, M_N\}$.
		The $N$th partial sum can be bounded above as follows:
		\begin{align*}
			\sum_{i=1}^N |\ell_i| &\le
			\sum_{i=1}^N |\ell_i - x_{n,i}| + \sum_{i=1}^N |x_{n,i}|.\\
			&< \sum_{i=1}^N 2^{-i} + \|\mb x_n\|\\
			&< 1 + \|\mb x_n\|.
		\end{align*}
		Additionally, since Cauchy sequences are bounded,
		we can bound $\|\mb x_n\|$ above by a constant $C$
		for all $n$, as, for an arbitrary fixed $r$,
		we can furnish a constant upper bound in the form
		\[\|\mb x_n\| \le \|\mb x_n - \mb x_r\| + \|\mb x_r\|
		\le \operatorname{diam}\{\mb x_n\} + \|\mb x_r\| =: C.\]
		Therefore we have bounded the partial sums of the $|\ell_i|$
		by the constant $1 + C$, so the infinite series converges.
	\end{proof}
	
	Now, we just need to prove that $\mb x_n\to \ell$
	in the metric of $\ell^1$. This is done with the same idea
	of bounding partial sums. Let $\eps>0$ and $N\in\NN$.
	For each $i\in\{1, 2, \dots, N\}$, we can find
	thresholds $M_i$ for which
	\[m > M_i \implies |\ell - x_{m,i}| < \eps\cdot 2^{-i - 1}.\]
	This means that, if $m > \max\{M_1, \dots, M_N\}$, then
	\[\sum_{i=1}^N |\ell_i - x_{m,i}| < \frac\eps2.\]
	We also know that $(\mb x_n)$ being Cauchy means that
	there is an $M$ such that
	\[m,\,n > M \implies \|\mb x_m - \mb x_n\| < \frac\eps2.\]
	Putting this together, we can fix
	$m$ to be larger than $M$ and all of the $M_i$
	and use the triangle inequality to say that
	$n > M$ implies
	\begin{align*}
		\sum_{i=1}^N |\ell_i - x_{n,i}| &\le
		\sum_{i=1}^N |\ell_i - x_{m,i}| + \sum_{i=1}^N |x_{m,i} - x_{n,i}|\\
		&\le \sum_{i=1}^N \eps\cdot 2^{-i-1} +
		\sum_{i=1}^\infty |x_{m,i} - x_{n,i}|\\
		&< \frac\eps2 + \frac\eps2 < \eps.
	\end{align*}
	In other words, as long as $n > M$ (noting that $M$ only depends on $\eps$),
	we can bound all the partial sums of $N$ below $\eps$,
	so $n > M$ implies that
	\[\|\mb x - \mb x_n\| < \eps,\]
	which is what it means for $\mb x_n$ to converge to $\mb x$
	in the metric of $\ell^1$, so we are done.
\end{soln}

\begin{problem}
	Prove that $f_n\to f$ and $g_n\to g$ uniformly on a set $E\subseteq\RR$
	then $f_n + g_n \to f + g$ uniformly on $E$.
	Show an example where $f_n\to f$ and $g_n\to g$ on some set
	$E\subseteq\RR$ but it is not true that
	$f_n\cdot g_n\to f\cdot g$ on $E\subseteq\RR$.
\end{problem}

\begin{soln}
	If $f_n\to f$ and $g_n\to g$ uniformly on $E\subseteq\RR$,
	then we can say that for each $\eps > 0$, there exists an $N$ such that
	$n > N$ implies
	\[|f(x) - f_n(x)| < \frac\eps2,\qquad |g(x) - g_n(x)| < \frac\eps2\]
	for all $x\in E$. For such $n$,
	\[|f(x) + g(x) - f_n(x) - g_n(x)| \le |f(x) - f_n(x)| + |g(x) - g_n(x)|
	< \eps.\]
	Again, noting that this holds for all $x$, we can conclude that
	$f_n + g_n\to f + g$ uniformly.\\
	
	For products:
	let $E = \RR$ and $f_n(x) = g_n(x) = x + \frac 1n$ for all $x$.
	Then, $f(x) = x$ and $g(x) = x$ are the uniform limits
	of $f_n$ and $g_n$ respectively because
	\[|f(x) - f_n(x)| = |g(x) - g_n(x)| = \frac 1n < \eps\]
	whenever $n > \frac 1\eps$, and the threshold for $n$ is independent of $x$.
	On the other hand,
	\[f_n(x)\cdot g_n(x) = \left(x + \frac 1n\right)^2
	= x^2 + \frac{2x}{n} + \frac{1}{n^2},\qquad
	f(x)g(x) = x^2.\]
	In this case, for a fixed $n$,
	\[|f(x)g(x) - f_n(x)g_n(x)| =
	\abs{\frac{2x}{n} + \frac{1}{n^2}}\]
	can be arbitrarily large, so any constant upper bound $\eps$
	cannot hold for all $x$; therefore, $f_ng_n$ does not
	converge to $fg$ uniformly.
\end{soln}

\begin{problem}
	Suppose that $f_n\to f$ uniformly on an interval $(a, b)$.
	Prove that if the $f_n$ are uniformly continuous functions
	then so is their uniform limit $f$.
\end{problem}

\begin{soln}
	Let the codomain of the $f_n$ be the metric space $(X, d)$
	and let $\eps > 0$.
	Note that, for any $x,\,y\in(a,b)$,
	\[d(f(x), f(y)) \le d(f(x), f_n(x)) + d(f_n(x), f_n(y)) + d(f_n(y), f(y)).\]
	Since $f_n\to f$ uniformly, there exists an $N$ such that
	$n > N$ implies
	\[d(f(z), f_n(z)) < \frac\eps3\]
	for all $z\in X$, so fix $n$ to be an integer greater than $N$.
	Now, $f_n$ is uniformly continuous, there exists a $\delta > 0$ such that
	\[|x - y| < \delta \implies d(f_n(x), f_n(y)) < \frac\eps3.\]
	Thus, for each $\eps > 0$, there exists a $\delta > 0$ for which
	\[|x - y| < \delta \implies d(f(x), f(y)) < \eps,\]
	which is what it means for $f$ to be uniformly continuous.
\end{soln}
















\end{document}